
\chapter{Trabalhos Relacionados}
\label{cap:trabalhos-relacionados}

Este capítulo apresenta e discute trabalhos relacionados à elicitação e especificação de requisitos de usabilidade no contexto da Engenharia de Software, com ênfase em métodos, técnicas e abordagens aplicadas em ambientes ágeis. O objetivo desta seção é posicionar o presente trabalho em relação à literatura existente, identificando convergências, limitações e lacunas que justificam a realização de uma meta-análise empírica do método USARP.

Os estudos analisados foram identificados a partir de buscas em bases reconhecidas, como Google Scholar, ACM Digital Library, IEEE Xplore e repositório de trabalhos acadêmicos da Universidade
Federal do Ceará, priorizando trabalhos que abordam usabilidade, requisitos não funcionais, \textit{personas}, \textit{user stories} e evidências empíricas de aplicação em contextos acadêmicos ou industriais.

\section{Métodos para Elicitação de Requisitos de Usabilidade}

Juristo, Moreno e Sánchez-Segura propuseram um conjunto de diretrizes para elicitação de funcionalidades de usabilidade, introduzindo o conceito de \textit{Functional Usability Features} (FUFs), que defende a incorporação de aspectos de usabilidade como funcionalidades explícitas do sistema \cite{juristo2007}. Embora o trabalho represente um avanço conceitual importante ao tratar a usabilidade como parte do comportamento funcional do software, sua abordagem é fortemente orientada a documentação tradicional de requisitos, o que dificulta sua adoção em ambientes ágeis baseados em artefatos leves.

Moreno e Yagüe exploraram a integração das FUFs com \textit{user stories}, investigando como mecanismos de usabilidade podem influenciar a criação, modificação ou refinamento de histórias de usuário e critérios de aceitação \cite{moreno2012}. Apesar dos resultados positivos, o estudo concentra-se predominantemente na fase de especificação, não oferecendo um processo estruturado para a elicitação sistemática dos requisitos de usabilidade a partir das características dos usuários.

Choma et al. propuseram a abordagem \textit{UserX Story}, que incorpora aspectos de Experiência do Usuário (UX) à escrita de \textit{user stories} por meio do uso de \textit{personas} e das heurísticas de Nielsen \cite{choma2016}. Nessa abordagem, as heurísticas são utilizadas principalmente como base para critérios de aceitação. Embora a proposta contribua para aproximar UX e métodos ágeis, ela trata a usabilidade como requisito não funcional e foca majoritariamente em aspectos de interface, não explorando seu impacto funcional ou arquitetural.

\section{Abordagens Baseadas em Personas e Prototipação}

Ferreira et al. propuseram a técnica PATHY, voltada à criação de \textit{personas} com foco na identificação de requisitos potenciais de software \cite{ferreira2018}. A técnica se mostrou eficaz para compreender o perfil dos usuários e suas necessidades, porém não é direcionada especificamente à elicitação de requisitos de usabilidade, exigindo sua combinação com outras abordagens para esse fim.

Outros estudos investigaram processos que combinam elicitação de requisitos e prototipação. Teixeira et al. apresentaram uma metodologia aplicada ao domínio da saúde, na qual protótipos são utilizados como mediadores para validação iterativa de requisitos \cite{teixeira2014}. Apesar de evidenciar benefícios práticos, o foco do estudo não está na usabilidade como requisito sistematicamente elicitado, mas na aceitação da solução pelos usuários finais.

Gonçalves e Rocha propuseram um processo metodológico baseado em normas internacionais para o desenvolvimento de interfaces inteligentes, integrando práticas de IHC ao ciclo de vida do software \cite{goncalves2019}. O trabalho destaca a importância da usabilidade, porém carece de evidências empíricas de aplicação contínua em contextos reais de desenvolvimento.

\section{O Método USARP e Evidências Empíricas}

O método USARP foi proposto por Oliveira Junior et al. com o objetivo de apoiar a elicitação e especificação de requisitos de usabilidade em fases iniciais do desenvolvimento, integrando \textit{personas}, \textit{user stories} e mecanismos de usabilidade derivados das FUFs \cite{oliveirajr2020}. Diferentemente de abordagens anteriores, o USARP oferece um processo estruturado que conecta explicitamente características dos usuários aos requisitos de usabilidade, promovendo discussões colaborativas entre membros da equipe.

Estudos subsequentes avaliaram o método em diferentes contextos, incluindo ambientes distribuídos, indicando benefícios relacionados à conscientização sobre usabilidade e à melhoria da qualidade das especificações produzidas \cite{marques2022}. Entretanto, esses estudos também evidenciam variações na forma como os mecanismos do método são utilizados, bem como diferenças entre percepção positiva e uso efetivo dos artefatos propostos.

\section{Síntese Crítica e Posicionamento do Trabalho}

A análise dos trabalhos relacionados evidencia que, embora existam diversas propostas para incorporar usabilidade ao processo de Engenharia de Requisitos, ainda são limitadas as abordagens que investigam simultaneamente a elicitação e a especificação de requisitos de usabilidade em contextos ágeis, com base em evidências empíricas sistematizadas.

Além disso, observa-se que a maioria dos estudos apresenta avaliações pontuais ou contextuais, sem uma síntese integrada das evidências disponíveis. Nesse sentido, o presente trabalho se diferencia ao realizar uma meta-análise empírica de braço único do método USARP, combinando análise estatística descritiva, redução de dimensionalidade por meio de PCA e análise qualitativa. Essa abordagem permite investigar padrões recorrentes de uso, percepção e adoção do método, contribuindo para uma compreensão mais abrangente de suas potencialidades e limitações em diferentes contextos de desenvolvimento de software.

\begin{table}[h!]
\captionsetup{width=15cm}
\caption{Comparação entre os trabalhos relacionados e o trabalho proposto}
\label{tab:comparacao-trabalhos}
\IBGEtab{}{
\begin{tabular}{lcccccc}
\toprule
\textbf{Trabalho} &
\textbf{Eng. de} &
\textbf{Elicitação} &
\textbf{Requisitos} &
\textbf{Personas /} &
\textbf{Evidência} &
\textbf{Síntese} \\
 & \textbf{Requisitos} &
\textbf{e Especificação} &
\textbf{de Usabilidade} &
\textbf{User Stories} &
\textbf{Empírica} &
\textbf{Sistematizada} \\
\midrule \midrule
Juristo et al. (2007)        & X & X & X &  & X &  \\
Moreno e Yagüe (2012)       & X &   & X & X & X &  \\
Choma et al. (2016)         & X &   & X & X & X &  \\
Ferreira et al. (2018)      & X & X &   & X & X &  \\
Oliveira Jr. et al. (2020)  & X & X & X & X & X &  \\
Marques et al. (2022)       & X &   & X & X & X &  \\
\textbf{Este trabalho}      & X & X & X & X & X & X \\
\bottomrule
\end{tabular}
}{
\Fonte{Elaborada pelo autor (2025).}
}
\end{table}
